\chapter[Equivalent Nodal Loads]{Equivalent Nodal Loads; Self-straining Problems}
\label{ch:eqload}

\section{Equivalent nodal loads}
The stiffness equation for a member subjected to intermediate concentrated loads or distributed loads can be written, with respect to the local or member axes, as
\begin{equation}
  \kmat\bs{\Delta}_l + \mat{f}_{fef} = \mat{f} \tag{a}\label{eq10:a}
\end{equation}
where $\mat{f}_{fef}$ is the fixed end force vector. Defining the equivalent force vector as $\mat{f}_{eq} = -\mat{f}_{fef}$, the above equation can be rewritten as
\begin{equation}
  \kmat\bs{\Delta}_l - \mat{f}_{eq} = \mat{f} \tag{b}\label{eq10:b}
\end{equation}
The above member stiffness equation is referred to the member (local) coordinate system. In order to carry out the assembling process, the stiffness equation referred to the structural (global) coordinate system is needed, which is:
\begin{equation}
  \Tmat^\tr\kmat\Tmat\bs{\Delta}_g - \Tmat^\tr\mat{f}_{eq} = \Tmat^\tr\mat{f}
  \quad\text{or}\quad
  \kmat_g\bs{\Delta}_g - \mat{f}_{eqg} = \mat{f}_g \tag{c}\label{eq10:c}
\end{equation}
where $\kmat_g = \Tmat^\tr\kmat\Tmat$, $\mat{f}_{eqg} = \Tmat^\tr\mat{f}_{eq}$, $\mat{f}_g = \Tmat^\tr\mat{f}$, and $\Tmat$ is the transformation matrix (see Ch.~\ref{ch:transform}). Eq.~\eqref{eq10:c} can be derived by using the contragredient principle and the congruent transformation. It can also be derived by premultiplying Eq.~\eqref{eq10:b} by $\Tmat^\tr$ and by using the relation that $\bs{\Delta}_l = \Tmat\bs{\Delta}_g$.

The stiffness equation for a structural system, obtained by assembling Eq.~\eqref{eq10:c}, takes the form
\begin{equation}
  \Kmat\mat{U} - \Pvec_{eq} = \F \tag{d}\label{eq10:d}
\end{equation}
where $\Kmat =$ assemblage of $\kmat_g$, $\mat{U} =$ collection of $\bs{\Delta}_g$, $\Pvec_{eq} =$ assemblage of $\Tmat^\tr\mat{f}_{eq}$, and $\F =$ assemblage of $\mat{f}_g$.

The \emph{nodal equilibrium equation} (considering the free body of a joint) states that
\begin{equation}
  \F = \Pvec_c \tag{e}\label{eq10:e}
\end{equation}
which means the nodal internal forces equal the external forces since $\Pvec_c =$ the applied nodal concentrated load vector. From Eqs.~\eqref{eq10:d} and \eqref{eq10:e}, we have
\begin{equation}
  \Kmat\mat{U} = \Pvec_c + \Pvec_{eq} \tag{f}\label{eq10:f}
\end{equation}
Once $\mat{U}$ is solved from Eq.~\eqref{eq10:f}, the member end forces, referred to the member coordinate system, can then be calculated according to the member stiffness equation, i.e.\ $\kmat\bs{\Delta}_l + \mat{f}_{fef} = \mat{f}$. To this end, we need to find $\bs{\Delta}_g$ first, which is a subset of $\mat{U}$, and then calculate $\bs{\Delta}_l$ by using $\bs{\Delta}_l = \Tmat\bs{\Delta}_g$.

From the above discussions, clearly the determination of $\mat{f}_{fef}$ or $\mat{f}_{eq}$ is the key element when dealing with members with intermediate loads or distributed loads. Table~5.1, on page 111 of the textbook, has listed fixed-end forces for some commonly encountered load types.
\begin{figure}[H]
  \centering
  % Fixed-end forces for commonly encountered load types (after Table 5.1 of the textbook)
\begin{tikzpicture}[scale=0.6,font=\scriptsize]
  \newcommand{\ffbeam}{\draw[member] (0,0) -- (6,0); \pic at (0,0) {wallL}; \pic at (6,0) {wallR};}
  \newcommand{\reactions}[4]{% #1 left moment label, #2 right moment label, #3 left shear, #4 right shear
    \draw[->,thick] (-0.9,0.5) arc (150:-30:0.5); \node[left] at (-1.0,0.55) {#1};
    \draw[->,thick] (6.9,0.5) arc (30:210:0.5); \node[right] at (7.0,0.55) {#2};
    \draw[load] (-0.6,-1.2) -- (-0.6,-0.1) node[pos=0,below] {#3};
    \draw[load] (6.6,-1.2) -- (6.6,-0.1) node[pos=0,below] {#4};}
  % ---------- (1) uniform load ----------
  \begin{scope}
    \ffbeam
    \foreach \x in {0.4,0.9,...,5.7}{\draw[->] (\x,0.9) -- (\x,0.1);} \draw (0.4,0.9) -- (5.7,0.9);
    \node[above] at (3,0.95) {$q$};
    \reactions{$\frac{qL^2}{12}$}{$\frac{qL^2}{12}$}{$\frac{qL}{2}$}{$\frac{qL}{2}$}
    \draw[dim] (0.6,-0.6) -- (5.4,-0.6) node[midway,fill=white,inner sep=1pt] {$L$};
  \end{scope}
  % ---------- (2) central point load ----------
  \begin{scope}[xshift=13cm]
    \ffbeam
    \draw[load] (3,1.3) -- (3,0.1) node[pos=0,right] {$P$};
    \reactions{$\frac{PL}{8}$}{$\frac{PL}{8}$}{$\frac P2$}{$\frac P2$}
    \draw[dim] (0.6,-0.6) -- (3,-0.6) node[midway,fill=white,inner sep=1pt] {$\frac L2$};
    \draw[dim] (3,-0.6) -- (5.4,-0.6) node[midway,fill=white,inner sep=1pt] {$\frac L2$};
    \draw (3,-0.8) -- (3,-0.4);
  \end{scope}
  % ---------- (3) point load at a ----------
  \begin{scope}[yshift=-4.2cm]
    \ffbeam
    \draw[load] (2,1.3) -- (2,0.1) node[pos=0,right] {$P$};
    \reactions{$\frac{Pab^2}{L^2}$}{$\frac{Pa^2b}{L^2}$}{$\frac{Pb^2(3a+b)}{L^3}$}{$\frac{Pa^2(a+3b)}{L^3}$}
    \draw[dim] (0.6,-0.6) -- (2,-0.6) node[midway,fill=white,inner sep=1pt] {$a$};
    \draw[dim] (2,-0.6) -- (5.4,-0.6) node[midway,fill=white,inner sep=1pt] {$b$};
    \draw (2,-0.8) -- (2,-0.4);
  \end{scope}
  % ---------- (4) support settlement ----------
  \begin{scope}[xshift=13cm,yshift=-4.2cm]
    \pic at (0,0) {wallL}; \pic at (6,-0.7) {wallR};
    \draw[member] (0,0) .. controls (2,0) and (4,-0.7) .. (6,-0.7);
    \draw[dim] (6.5,0) -- (6.5,-0.7) node[midway,right] {$\Delta$};
    \draw[->,thick] (-0.9,0.5) arc (150:-30:0.5); \node[left] at (-1.0,0.55) {$\frac{6EI\Delta}{L^2}$};
    \draw[->,thick] (6.9,-0.2) arc (30:-150:0.5); \node[right] at (7.0,0.2) {$\frac{6EI\Delta}{L^2}$};
    \draw[load] (-0.6,-1.2) -- (-0.6,-0.1) node[pos=0,below] {$\frac{12EI\Delta}{L^3}$};
    \draw[load] (6.6,-0.8) -- (6.6,-1.9) node[below] {$\frac{12EI\Delta}{L^3}$};
    \draw[dim] (0.6,-1.4) -- (5.4,-1.4) node[midway,fill=white,inner sep=1pt] {$L$};
  \end{scope}
  % ---------- (5) concentrated moment ----------
  \begin{scope}[yshift=-8.4cm]
    \ffbeam
    \draw[->,thick] (2.4,0.5) arc (100:-100:0.5); \node[right] at (2.6,-0.35) {$M$};
    \reactions{$\frac{Mb(2a-b)}{L^2}$}{$\frac{Ma(2b-a)}{L^2}$}{$\frac{6Mab}{L^3}$}{$\frac{6Mab}{L^3}$}
    \draw[dim] (0.6,-0.6) -- (2.2,-0.6) node[midway,fill=white,inner sep=1pt] {$a$};
    \draw[dim] (2.2,-0.6) -- (5.4,-0.6) node[midway,fill=white,inner sep=1pt] {$b$};
    \draw (2.2,-0.8) -- (2.2,-0.4);
  \end{scope}
  % ---------- (6) triangular load ----------
  \begin{scope}[xshift=13cm,yshift=-8.4cm]
    \ffbeam
    \foreach \x in {0.9,1.4,...,5.7}{\draw[->] (\x,{0.1+0.16*\x}) -- (\x,0.1);} \draw (0.4,0.1) -- (5.7,1.0);
    \node[above] at (4.5,1.0) {$q$};
    \reactions{$\frac{qL^2}{30}$}{$\frac{qL^2}{20}$}{$\frac{3qL}{20}$}{$\frac{7qL}{20}$}
    \draw[dim] (0.6,-0.6) -- (5.4,-0.6) node[midway,fill=white,inner sep=1pt] {$L$};
  \end{scope}
  % ---------- (7) symmetric triangular load ----------
  \begin{scope}[yshift=-12.6cm]
    \ffbeam
    \foreach \x in {0.7,1.2,...,2.9}{\draw[->] (\x,{0.1+0.34*\x}) -- (\x,0.1);}
    \foreach \x in {3.4,3.9,...,5.5}{\draw[->] (\x,{0.1+0.34*(6-\x)}) -- (\x,0.1);}
    \draw (0.3,0.1) -- (3,1.1) -- (5.7,0.1); \node[above] at (3,1.1) {$q$};
    \reactions{$\frac{5qL^2}{96}$}{$\frac{5qL^2}{96}$}{$\frac{qL}{4}$}{$\frac{qL}{4}$}
    \draw[dim] (0.6,-0.6) -- (5.4,-0.6) node[midway,fill=white,inner sep=1pt] {$L$};
  \end{scope}
  % ---------- (8) partial uniform load ----------
  \begin{scope}[xshift=13cm,yshift=-12.6cm]
    \ffbeam
    \foreach \x in {1.8,2.15,...,4.2}{\draw[->] (\x,0.8) -- (\x,0.1);} \draw (1.8,0.8) -- (4.2,0.8);
    \node[above right] at (4.2,0.8) {$q$};
    \draw[dim] (1.8,1.2) -- (3,1.2) node[midway,fill=white,inner sep=1pt] {$\frac d2$};
    \draw[dim] (3,1.2) -- (4.2,1.2) node[midway,fill=white,inner sep=1pt] {$\frac d2$};
    \draw[->,thick] (-0.9,0.5) arc (150:-30:0.5); \node[left,align=right] at (-1.0,0.55) {$\frac{qd}{L^2}\cdot$\\[-2pt]$[ab^2+\frac{(a-2b)d^2}{12}]$};
    \draw[->,thick] (6.9,0.5) arc (30:210:0.5); \node[right,align=left] at (7.0,0.55) {$\frac{qd}{L^2}\cdot$\\[-2pt]$[a^2b+\frac{(b-2a)d^2}{12}]$};
    \draw[load] (-0.6,-1.2) -- (-0.6,-0.1) node[pos=0,below] {$R_1$};
    \draw[load] (6.6,-1.2) -- (6.6,-0.1) node[pos=0,below] {$R_2$};
    \draw[dim] (0.6,-0.6) -- (3,-0.6) node[midway,fill=white,inner sep=1pt] {$a$};
    \draw[dim] (3,-0.6) -- (5.4,-0.6) node[midway,fill=white,inner sep=1pt] {$b$};
    \draw (3,-0.8) -- (3,-0.4);
    \node[anchor=north west,align=left] at (-1.5,-2.0) {$R_1 = \frac{qd}{L^3}\bigl[(2a+L)b^2+\bigl(\frac{a-b}{4}\bigr)d^2\bigr]$\\[4pt]$R_2 = \frac{qd}{L^3}\bigl[(2b+L)a^2-\bigl(\frac{a-b}{4}\bigr)d^2\bigr]$};
  \end{scope}
  \node[anchor=north west,align=left,font=\footnotesize] at (-1.5,-14.6) {Sign convention: all loads and reactive forces\\are positive in the directions shown.};
\end{tikzpicture}

  \caption{Fixed-end forces for common load types (after Table~5.1 of the textbook).}
  \label{fig:fef-table}
\end{figure}

In general, we can determine $\mat{f}_{eq}$ using the principle of virtual displacements. In this way, we let the virtual work done by $\mat{f}_{eq}$ be equal to that done by the intermediate concentrated loads (can be expressed as delta functions) or distributed loads. That is,
\begin{gather}
  \delta\bs{\Delta}_l^\tr\mat{f}_{eq} = \int_0^L \delta v(x)^\tr q(x)\,dx = \int_0^L (\NN\delta\bs{\Delta}_l)^\tr q(x)\,dx
  \;\rightarrow\; \delta\bs{\Delta}_l^\tr\mat{f}_{eq} = \delta\bs{\Delta}_l^\tr\int_0^L \NN^\tr q(x)\,dx\\
  \rightarrow\quad \mat{f}_{eq} = \int_0^L \NN^\tr q(x)\,dx
\end{gather}
where $\NN = \langle N_1, N_2, \dots\rangle$ are shape functions. Written in component form,
\begin{equation}
  (\mat{f}_{eq})_i = \int_0^L N_i\,q(x)\,dx
\end{equation}
Examples illustrating the above equation (will be done in class).

\begin{example}
Consider the beam shown below with the properties: $E = 28$~GPa, $I = 5.8\times10^9$~mm$^4$.
\begin{fig}
  % Example beam: 4 nodes, 3 elements, with loads
\begin{tikzpicture}[scale=0.7,font=\small]
  \draw[member] (0,0) -- (20,0);
  \pic at (0,0) {wallL}; \pic at (10,0) {roller}; \pic at (20,0) {roller};
  \foreach \x/\n in {0/1,6/2,10/3,20/4}{\node[below right,inner sep=2pt] at (\x,0) {\n};}
  \node[node dot] at (6,0) {};
  % triangular load on element 1
  \foreach \x in {0.5,1.2,...,5.6}{\draw[->] (\x,{0.2+0.28*\x}) -- (\x,0.15);} \draw (0,0.2) -- (6,1.9);
  \node[above left] at (1.5,1.6) {30\,kN/m};
  \draw[load] (6,2.8) -- (6,0.2) node[pos=0,above] {200\,kN};
  \draw[->,thick] (11.2,1.3) arc (30:180:0.9); \node[above] at (10.2,1.6) {90\,kNm};
  \draw[load] (15,2.4) -- (15,0.2) node[pos=0,above] {150\,kN};
  \node[below] at (3,-0.4) {$1.5I$}; \node[below] at (15,-0.4) {$I$};
  \draw[dim] (0,-1.8) -- (6,-1.8) node[midway,fill=white,inner sep=1pt] {6\,m};
  \draw[dim] (6,-1.8) -- (10,-1.8) node[midway,fill=white,inner sep=1pt] {4\,m};
  \draw[dim] (10,-1.8) -- (15,-1.8) node[midway,fill=white,inner sep=1pt] {5\,m};
  \draw[dim] (15,-1.8) -- (20,-1.8) node[midway,fill=white,inner sep=1pt] {5\,m};
  \foreach \x in {0,6,10,15,20}{\draw (\x,-1.5) -- (\x,-2.1);}
\end{tikzpicture}

\end{fig}
\begin{fig}
  % Fixed-end forces of elements 1 and 3
\begin{tikzpicture}[scale=0.7,font=\small]
  % element 1
  \begin{scope}
    \draw[member] (0,0) -- (6,0); \pic at (0,0) {wallL}; \pic at (6,0) {wallR};
    \foreach \x in {0.5,1.2,...,5.6}{\draw[->] (\x,{0.2+0.28*\x}) -- (\x,0.15);} \draw (0,0.2) -- (6,1.9);
    \node[above left] at (1.2,1.5) {30\,kN/m};
    \node[below right] at (0,0) {1}; \node[below left] at (6,0) {2};
    \draw[->,thick] (-0.5,-0.6) arc (-100:100:0.6); \node[above left] at (-0.4,0.5) {54};
    \draw[->,thick] (6.5,-0.6) arc (280:80:0.6); \node[above right] at (6.4,0.5) {$-36$};
    \draw[load] (-0.5,-1.4) -- (-0.5,-0.1) node[pos=0,below] {63};
    \draw[load] (6.5,-1.4) -- (6.5,-0.1) node[pos=0,below] {27};
  \end{scope}
  % element 3
  \begin{scope}[xshift=10.5cm]
    \draw[member] (0,0) -- (10,0); \pic at (0,0) {wallL}; \pic at (10,0) {wallR};
    \draw[load] (5,2.2) -- (5,0.15) node[pos=0,above] {150\,kN};
    \node[below right] at (0,0) {3}; \node[below left] at (10,0) {4};
    \draw[->,thick] (-0.5,-0.6) arc (-100:100:0.6); \node[above left] at (-0.4,0.5) {187.5};
    \draw[->,thick] (10.5,-0.6) arc (280:80:0.6); \node[above right] at (10.4,0.5) {$-187.5$};
    \draw[load] (-0.5,-1.4) -- (-0.5,-0.1) node[pos=0,below] {75};
    \draw[load] (10.5,-1.4) -- (10.5,-0.1) node[pos=0,below] {75};
    \node[right] at (11.6,-0.6) {$\mat{f}_{fef}$};
  \end{scope}
\end{tikzpicture}

\end{fig}
\begin{fig}
  % Equivalent nodal loads (solid) and concentrated nodal loads (dashed)
\begin{tikzpicture}[scale=0.62,font=\small]
  \draw[member] (0,0) -- (20,0);
  \foreach \x/\n in {0/1,6/2,10/3,20/4}{\node[below,inner sep=3pt] at (\x,0) {\n};}
  % node 1
  \draw[load] (0,0.2) -- (0,1.6) node[above] {$-63$};
  \draw[->,thick] (0.6,0.4) arc (30:180:0.6); \node[left] at (-0.5,0.4) {$-54$};
  % node 2
  \draw[load] (6,0.2) -- (6,1.6) node[above right] {$-27$};
  \draw[load,dashed] (6,1.7) -- (6,3.2) node[above] {$-200$};
  \draw[->,thick] (6.6,0.4) arc (30:180:0.6); \node[left] at (5.5,0.6) {36};
  % node 3
  \draw[load] (10,0.2) -- (10,1.6) node[above] {$-75$};
  \draw[->,thick] (10.6,0.4) arc (30:180:0.6); \node[right] at (10.6,0.7) {$-187.5$};
  \draw[->,thick,dashed] (10.9,0.9) arc (20:180:1.0); \node[right] at (11.0,1.4) {$-90$};
  % node 4
  \draw[load] (20,0.2) -- (20,1.6) node[above] {$-75$};
  \draw[->,thick] (20.6,0.4) arc (30:180:0.6); \node[right] at (20.6,0.7) {187.5};
  % legend
  \draw (22.5,1.6) -- (24,1.6) node[right] {$\mat{f}_{eq}$};
  \draw[dashed] (22.5,0.6) -- (24,0.6) node[right] {$\Pvec_c$};
\end{tikzpicture}

\end{fig}
Note that in the program, the fixed end forces \texttt{FEF} shall be input as follows:
\begin{lstlisting}[language={},numbers=none,basicstyle=\small\ttfamily]
*FEF
1   63    54    27   -36
2    0     0     0     0
3   75   187.5  75  -187.5
\end{lstlisting}
After assembling, we obtain the final nodal loads:
\begin{fig}
  % Final nodal loads after assembling
\begin{tikzpicture}[scale=0.7,font=\small]
  \draw[member] (0,0) -- (20,0);
  \pic at (0,0) {wallL}; \pic at (10,0) {roller}; \pic at (20,0) {roller};
  \foreach \x/\n in {0/1,6/2,10/3,20/4}{\node[below right,inner sep=3pt] at (\x,0) {\n};}
  \draw[load] (6,0.2) -- (6,2.2) node[above] {$-227$\,kN};
  \draw[->,thick] (6.6,0.5) arc (30:180:0.7); \node[left] at (5.6,0.9) {36\,kNm};
  \draw[->,thick] (10.8,0.5) arc (30:180:0.9); \node[right] at (10.6,1.1) {$-277.5$\,kNm};
  \draw[->,thick] (20.8,0.5) arc (30:180:0.9); \node[above] at (20,1.3) {187.5\,kNm};
\end{tikzpicture}

\end{fig}
\begin{equation*}
  \Longrightarrow\quad \Pvec = \Pvec_c + \Pvec_{eq} = \Bmat{-200\\ 0\\ -90\\ 0} + \Bmat{-27\\ 36\\ -187.5\\ 187.5}
  = \Bmat{-227\text{ kN}\\ 36\text{ kN}\cdot\text{m}\\ -277.5\text{ kN}\cdot\text{m}\\ 187.5\text{ kN}\cdot\text{m}}
\end{equation*}
\begin{equation*}
  \Kmat = \bmat{43983 & 20300 & 60900 & 0\\ 20300 & 324800 & 81200 & 0\\ 60900 & 81200 & 227360 & 32480\\ 0 & 0 & 32480 & 64960}\quad\text{(with appropriate units)}
\end{equation*}
Then,
\begin{equation*}
  \bs{\Delta} = \Bmat{-4.4729\text{ m}\\ 0.56143\text{ rad}\\ -0.68415\text{ rad}\\ 3.2285\text{ rad}}\times10^{-3}
\end{equation*}

\topic{Force recovery (element one)}
\begin{fig}
  % Element 1: fixed-end forces and local displacements
\begin{tikzpicture}[scale=0.8,font=\small]
  \begin{scope}
    \draw[member] (0,0) -- (6,0);
    \node[below right] at (0.2,0) {1}; \node[below left] at (5.8,0) {2};
    \draw[->,thick] (-0.15,-0.6) arc (-100:100:0.6); \node[above left] at (-0.1,0.5) {54};
    \draw[->,thick] (6.15,-0.6) arc (280:80:0.6); \node[above right] at (6.1,0.5) {$-36$};
    \draw[load] (0,-1.2) -- (0,-0.1) node[pos=0,below] {63};
    \draw[load] (6,-1.2) -- (6,-0.1) node[pos=0,below] {27};
  \end{scope}
  \begin{scope}[xshift=9.5cm]
    \draw[member] (0,0) -- (6,0);
    \node[below right] at (0.2,0) {1}; \node[below left] at (5.8,0) {2};
    \draw[->,thick] (-0.15,-0.6) arc (-100:100:0.6); \node[above left] at (-0.1,0.5) {0};
    \draw[->,thick] (6.15,-0.6) arc (280:80:0.6); \node[above right] at (6.1,0.5) {$\Delta_2$};
    \draw[load] (0,-1.2) -- (0,-0.1) node[pos=0,below] {0};
    \draw[load] (6,-1.2) -- (6,-0.1) node[pos=0,below] {$\Delta_1$};
  \end{scope}
\end{tikzpicture}

\end{fig}
Fixed end forces $\mat{f}_{fef} = \Bmat{63\\ 54\\ 27\\ -36}$, and local displacement vector $\bs{\Delta}_l = \Bmat{0\\ 0\\ \Delta_1\\ \Delta_2}$. Then from Eq.~\eqref{eq10:a},
\begin{equation*}
  \mat{f}_1 = \kmat\bs{\Delta}_l + \mat{f}_{fef} = \Bmat{146.3\text{ kN}\\ 281.19\text{ kN}\cdot\text{m}\\ -56.33\text{ kN}\\ 236.78\text{ kN}\cdot\text{m}}.
\end{equation*}
Similarly, we can obtain the member forces of elements 2 and 3 as follows:
\begin{equation*}
  \mat{f}_2 = \Bmat{-143.67\text{ kN}\\ -236.78\text{ kN}\cdot\text{m}\\ 143.67\text{ kN}\\ -337.92\text{ kN}\cdot\text{m}}
  \quad\text{and}\quad
  \mat{f}_3 = \Bmat{99.79\text{ kN}\\ 247.92\text{ kN}\cdot\text{m}\\ 50.21\text{ kN}\\ 0\text{ kN}\cdot\text{m}}
\end{equation*}
\end{example}

\section{Self-straining problems}
Self-straining problems include those dealing with temperature changes, fabrication errors (misfit), and support settlement. For these kinds of problems, the member stiffness equations can also be written as
\begin{equation}
  \kmat\bs{\Delta}_l + \mat{f}_{fef} = \mat{f}\qquad\text{or}\qquad \kmat\bs{\Delta}_l = \mat{f} + \mat{f}_{eq}
\end{equation}
where $\mat{f}_{eq} = -\mat{f}_{fef}$ are defined as the forces that need to be applied to the member so that all nodal displacements are equal to zero, or equivalently defined as the forces that counterbalance the initial strains. For example, if an axial member is subjected to a temperature rise $\Delta T$, then obviously a compressive force with the magnitude of $AE\alpha\Delta T$ needs to be applied to the ends of the member in order not to induce deformations or forces for the other members, where $\alpha =$ thermal expansion coefficient, a material constant. Actually, this compressive force produces an axial strain $= -\alpha\Delta T$, which counterbalances the initial strain due to the temperature rise $= \alpha\Delta T$. In summary, if $\eps_i = \alpha\Delta T$ represents the initial strain, then the fixed-force $f_{fef} = -EA\eps_i$ and therefore the equivalent force $f_{eq} = -f_{fef} = EA\eps_i$.

Similarly, if an axial bar has the fabrication error $\Delta L$, then it has the initial strain $\eps_i = \Delta L/L$. The fixed-end force $f_{fef} = -EA\eps_i$, and the equivalent force $f_{eq} = EA\eps_i = EA\Delta L/L$.

In general, from the principle of virtual work, we have, for an axial bar, that
\begin{equation}
  f_{eq} = EA\frac1L\int_0^L \eps_i(x)\,dx
\end{equation}
For beams, the fixed-end forces corresponding to support settlement or temperature changes can be determined similarly as we discuss in class.
\begin{fig}
  % Self-straining: temperature rise in two bars of a three-bar truss = locked + unlocked
\begin{tikzpicture}[scale=0.8,font=\small]
  \newcommand{\truss}{%
    \draw[member] (-2,0) -- (0,3) -- (2,0); \draw[member] (0,0) -- (0,3);
    \pic at (-2,0) {pin}; \pic at (0,0) {pin}; \pic at (2,0) {pin}; \node[node dot] at (0,3) {};}
  \begin{scope}
    \truss
    \node[above left] at (-1.2,1.8) {$\uparrow\Delta T_1$}; \node[above right] at (1.2,1.8) {$\downarrow\Delta T_2$};
    \node at (-1.2,3.9) {$\Delta T_1 > 0$}; \node at (1.2,3.9) {$\Delta T_2 > 0$};
  \end{scope}
  \node at (3.6,1.5) {$=$};
  \begin{scope}[xshift=7.2cm]
    \truss \node at (0,4.6) {\itshape locked};
    \draw[load] ($(0,3)+(-1.4,2.1)$) -- ($(0,3)+(-0.15,0.22)$) node[pos=0,above left] {$EA\alpha\Delta T_2$};
    \draw[load] ($(0,3)+(1.4,2.1)$) -- ($(0,3)+(0.15,0.22)$) node[pos=0,above right] {$EA\alpha\Delta T_1$};
  \end{scope}
  \node at (10.8,1.5) {$+$};
  \begin{scope}[xshift=14.4cm]
    \truss \node at (0,4.6) {\itshape unlocked};
    \draw[load] ($(0,3)+(-0.15,0.22)$) -- ($(0,3)+(-1.4,2.1)$) node[above left] {$EA\alpha\Delta T_2$};
    \draw[load] ($(0,3)+(0.15,0.22)$) -- ($(0,3)+(1.4,2.1)$) node[above right] {$EA\alpha\Delta T_1$};
  \end{scope}
\end{tikzpicture}

\end{fig}

\begin{example}[Example 5.10 in the textbook]
Member $ad$ of the truss is 3~mm fabricated too short and $E = 200{,}000$~MPa.
\begin{fig}
  % Example 5.10: truss with member ad fabricated 3 mm too short
\begin{tikzpicture}[scale=0.75,font=\scriptsize]
  \coordinate (c) at (0,0); \coordinate (b) at (10.928,0); \coordinate (a) at (6.928,4); \coordinate (d) at (0,4);
  \draw[member] (c) -- (a) -- (b) -- (c); \draw[member] (d) -- (a);
  \pic[rotate=-90] at (d) {pin}; \pic at (c) {pin}; \pic at (b) {roller};
  \node[above right] at (a) {$a$}; \node[below right] at (b) {$b$}; \node[below left] at (c) {$c$}; \node[above right] at (d) {$d$};
  \node[above,sloped] at ($(c)!0.5!(a)$) {$A = 15\times10^3$ mm$^2$};
  \node[above,sloped] at ($(a)!0.5!(b)$) {$A = 20\times10^3$ mm$^2$};
  \node[above] at (3.5,4) {$A = 20\times10^3$ mm$^2$};
  \node[below] at (5.5,0) {$A = 18\times10^3$ mm$^2$};
  \draw[->] (1.6,0) arc (0:30:1.6); \node at (2.0,0.4) {$30^\circ$};
  \draw[->] (9.4,0) arc (180:135:1.5); \node at (8.9,0.5) {$45^\circ$};
  \draw[dim] (0,4.8) -- (6.928,4.8) node[midway,fill=white,inner sep=1pt] {6.928 m};
  \draw[dim] (0,-1.2) -- (10.928,-1.2) node[midway,fill=white,inner sep=1pt] {10.928 m};
  \draw[dim] (11.8,0) -- (11.8,4) node[midway,right] {4 m};
  \draw[->] (-0.6,0.3) -- (-0.6,1.3) node[above] {$y$}; \draw[->] (11.2,-0.5) -- (12.2,-0.5) node[right] {$x$};
\end{tikzpicture}

\end{fig}
DOFs and fixed end forces:
\begin{fig}
  % Example 5.10: DOFs and reactions
\begin{tikzpicture}[scale=0.55,font=\small]
  \coordinate (c) at (0,0); \coordinate (b) at (10.928,0); \coordinate (a) at (6.928,4); \coordinate (d) at (0,4);
  \draw[thinmember] (c) -- (a) -- (b) -- (c); \draw[thinmember] (d) -- (a);
  \foreach \p in {a,b,c,d}{\fill (\p) circle (0.07);}
  \node[above left] at (a) {$a$}; \node[above] at (b) {$b$}; \node[above right] at (c) {$c$}; \node[below right] at (d) {$d$};
  \draw[load] (a) -- ++(1.6,0) node[right] {$P_1,\Delta_1$};
  \draw[load] (a) -- ++(0,1.8) node[right] {$P_2,\Delta_2$};
  \draw[load] (b) -- ++(1.6,0) node[right] {$P_3,\Delta_3$};
  \draw[load] ($(d)+(-1.8,0)$) -- (d) node[pos=0,below] {$R_{xd}$};
  \draw[load] ($(c)+(-1.8,0)$) -- (c) node[pos=0,below] {$R_{xc}$};
  \draw[load] ($(c)+(0,-1.8)$) -- (c) node[pos=0,right] {$R_{yc}$};
  \draw[load] ($(b)+(0,-1.8)$) -- (b) node[pos=0,right] {$R_{yb}$};
\end{tikzpicture}

\end{fig}
\begin{fig}
  % Example 5.10: fixed-end forces of member ad (3 mm too short)
\begin{tikzpicture}[scale=0.8,font=\small]
  \draw[member] (0,0) -- (6.928,0);
  \draw (6.928,-0.4) -- (6.928,0.9); \draw (7.3,-0.4) -- (7.3,0.9);
  \draw[dim] (7.3,0.7) -- (8.6,0.7); \node[right] at (8.6,0.7) {3 mm};
  \draw[<-] (6.928,0.7) -- (5.8,0.7);
  \node[below right] at (0,0) {$d$}; \node[above left] at (6.928,0) {$a$};
  \draw[load] (-1.6,0) -- (-0.1,0) node[pos=0,below] {$F^F_{xd}$};
  \draw[load] (8.6,0) -- (7.4,0) node[pos=0,right] {$F^F_{xa}$};
  \draw[dim] (0,-0.8) -- (6.928,-0.8) node[midway,fill=white,inner sep=1pt] {6.928 m};
\end{tikzpicture}

\end{fig}
\begin{equation*}
  F^F_{xa} = -F^F_{xd} = \Bigl(\frac{EA}{L}\Bigr)_{ad}\cdot\Delta L_{ad} = 577.37\times3 = 1732.1\text{ kN}
\end{equation*}
Obviously, the fixed end force is zero except member $ad$.

Q: $\mat{f}^{ad}_{fef} = \langle -1732.1, 1732.1\rangle$~kN or $\mat{f}^{ad}_{fef} = \langle 1732.1, -1732.1\rangle$~kN? Why?

Load vector:
\begin{equation*}
  \Pvec = \Pvec_c + \Pvec_{eq} = \Bmat{0\\ 0\\ 0} + \Bmat{-1732.1\\ 0\\ 0}\text{ kN}
\end{equation*}
Stiffness equations:
\begin{equation*}
  10^2\bmat{12.122 & -1.912 & -3.536\\ & 4.473 & 3.536\\ & & 6.830}\Bmat{\Delta_1\\ \Delta_2\\ \Delta_3} = \Bmat{-1732.1\\ 0\\ 0}
\end{equation*}
Displacement: $\langle\Delta_1, \Delta_2, \Delta_3\rangle = \langle -1.683, -0.0158, -8.448\rangle$~mm.
\end{example}

\begin{example}[Example 5.11 in the textbook]
The support at $b$ settles 15~mm; analyze the problem.
\begin{fig}
  % Example 5.11: two-span beam, support b settles 15 mm
\begin{tikzpicture}[scale=0.8,font=\small]
  \draw[member] (0,0) -- (13,0);
  \pic at (0,0) {pin}; \pic at (8,0) {roller}; \pic at (13,0) {wallR};
  \node[above] at (0,0.1) {$a$}; \node[above] at (8,0.1) {$b$}; \node[above left] at (13,0.1) {$c$};
  \draw[dim] (0,-1.0) -- (8,-1.0) node[midway,fill=white,inner sep=1pt] {8 m};
  \draw[dim] (8,-1.0) -- (13,-1.0) node[midway,fill=white,inner sep=1pt] {5 m};
  \foreach \x in {0,8,13}{\draw (\x,-0.8) -- (\x,-1.2);}
\end{tikzpicture}

\end{fig}
Fixed end forces (note that $\theta_{zb} = 0$ in the figure below):
\begin{fig}
  % Example 5.11: fixed-end state with support b settled 15 mm (theta_zb = 0)
\begin{tikzpicture}[scale=0.8,font=\small]
  \pic at (0,0) {wallL}; \pic at (13,0) {wallR};
  \draw[dashed] (0,0) -- (13,0);
  \draw[member] (0,0) .. controls (3,0) and (5,-0.9) .. (8,-0.9) .. controls (10,-0.9) and (11.5,0) .. (13,0);
  \pic at (8,-0.9) {roller};
  \draw[dim] (13.6,0) -- (13.6,-0.9) node[midway,right] {15 mm};
  \draw[dim] (8.4,0.6) -- (8.4,-0.9); \node[right] at (8.4,0.5) {15 mm};
  \node[above] at (0,0.1) {$a$}; \node[above right] at (8,-0.8) {$b$}; \node[above left] at (13,0.1) {$c$};
\end{tikzpicture}

\end{fig}
\begin{fig}
  % Example 5.11: fixed-end forces of members ab and bc
\begin{tikzpicture}[scale=0.8,font=\small]
  \begin{scope}
    \draw[member] (0,0) -- (7,0); \fill (0,0) circle (0.06); \fill (7,0) circle (0.06);
    \node[above right] at (0,0) {$a$}; \node[above left] at (7,0) {$b$};
    \draw[->,thick] (-0.1,-0.6) arc (-100:100:0.6); \node[above left] at (-0.1,0.5) {$M^F_{za}$};
    \draw[->,thick] (7.1,-0.6) arc (280:80:0.6); \node[above right] at (7.1,0.5) {$M^F_{zb}$};
    \draw[load] (0,-1.3) -- (0,-0.1) node[pos=0,below] {$F^F_{ya}$};
    \draw[load] (7,-1.3) -- (7,-0.1) node[pos=0,below] {$F^F_{yb}$};
  \end{scope}
  \begin{scope}[xshift=10cm]
    \draw[member] (0,0) -- (6,0); \fill (0,0) circle (0.06); \fill (6,0) circle (0.06);
    \node[above right] at (0,0) {$b$}; \node[above left] at (6,0) {$c$};
    \draw[->,thick] (-0.1,-0.6) arc (-100:100:0.6); \node[above left] at (-0.1,0.5) {$M^F_{zb}$};
    \draw[->,thick] (6.1,-0.6) arc (280:80:0.6); \node[above right] at (6.1,0.5) {$M^F_{zc}$};
    \draw[load] (0,-1.3) -- (0,-0.1) node[pos=0,below] {$F^F_{yb}$};
    \draw[load] (6,-1.3) -- (6,-0.1) node[pos=0,below] {$F^F_{yc}$};
  \end{scope}
\end{tikzpicture}

\end{fig}
\begin{minipage}{0.5\linewidth}
\begin{align*}
  M^F_{za} = M^F_{zb} &= 200(-18.75)(-15)\\ &= 56.25\times10^3\text{ kNmm}\\
  F^F_{ya} = -F^F_{yb} &= 200(-0.00469)(-15)\\ &= 14.07\text{ kN}
\end{align*}
\end{minipage}\hfill
\begin{minipage}{0.48\linewidth}
\begin{align*}
  M^F_{zb} = M^F_{zc} &= 200(12.00)(-15)\\ &= -36.0\times10^3\text{ kNmm}\\
  F^F_{yb} = -F^F_{yc} &= 200(0.00480)(-15)\\ &= -14.40\text{ kN}
\end{align*}
\end{minipage}
\begin{equation*}
  \mat{f}^{ab}_{fef} = \langle 14.07, 56250, -14.07, 56250\rangle^\tr;\qquad
  \mat{f}^{bc}_{fef} = \langle -14.4, -36000, 14.4, -36000\rangle^\tr
\end{equation*}
Load vector:
\begin{equation*}
  \Pvec = \Pvec_c + \Pvec_{eq} = \Bmat{0\\ 0} + \Bmat{-56.25\times10^3\\ -56.25\times10^3 + 36\times10^3}\text{ kN}\cdot\text{mm}
  = \Bmat{-56.25\times10^3\\ -20.25\times10^3}\text{ kN}\cdot\text{mm}
\end{equation*}
Stiffness equations:
\begin{equation*}
  200\times10^5\bmat{1 & 0.5\\ 0.5 & 1.4}\Bmat{\theta_{za}\\ \theta_{zb}} = \Bmat{-56.25\times10^3\\ -20.25\times10^3}
\end{equation*}
Displacement: $\langle\theta_{za}, \theta_{zb}\rangle = \langle -0.002984, 0.000342\rangle$~rad.

\topic{Comparison: the slope-deflection method}
\begin{fig}
  % Comparison with the slope-deflection method
\begin{tikzpicture}[scale=0.8,font=\small]
  \begin{scope}
    \draw[member] (0,0) -- (10,0); \pic at (0,0) {pin}; \pic at (4,0) {roller}; \pic at (10,0) {wallR};
    \draw[->,thick] (0.6,0.5) arc (20:170:0.6); \node[above] at (-0.1,0.9) {$\theta_a$};
    \draw[->,thick] (4.6,0.5) arc (20:170:0.6); \node[above] at (3.9,0.9) {$\theta_b$};
    \node[below] at (2,-0.6) {$EI_1, L_1$}; \node[below] at (7,-0.6) {$EI_2, L_2$};
  \end{scope}
  \begin{scope}[yshift=-3.6cm]
    \draw[member] (0,0) -- (10,0); \pic at (0,0) {pin}; \pic at (4,0) {roller}; \pic at (10,0) {wallR};
    \node[above] at (0,0.6) {$M_a = 0$};
    \draw[->,thick] (4.6,0.5) arc (20:170:0.6); \node[above] at (3.9,0.9) {$\theta_b$};
    \node[below] at (2,-0.6) {$EI_1, L_1$}; \node[below] at (7,-0.6) {$EI_2, L_2$};
  \end{scope}
\end{tikzpicture}

\end{fig}
\begin{equation*}
  \Kmat = \begin{bmatrix}\dfrac{4E_1I_1}{L_1} & \dfrac{2E_1I_1}{L_1}\\[8pt] \dfrac{2E_1I_1}{L_1} & \dfrac{4E_1I_1}{L_1}+\dfrac{4E_2I_2}{L_2}\end{bmatrix}
  \qquad\qquad
  k = \frac{3E_1I_1}{L_1} + \frac{4E_2I_2}{L_2},\quad M_b = k\cdot\theta_b
\end{equation*}
\end{example}

\section{Thermal load problem of beam}
\begin{fig}
  % Thermal load problem of a beam (after the textbook)
\begin{tikzpicture}[scale=0.8,font=\scriptsize]
  % ---- (a) beam with temperature gradient ----
  \begin{scope}
    \draw (0,0) rectangle (7,0.9);
    \pic at (0,0) {pin}; \pic at (7,0) {roller};
    \node[left] at (0,0.45) {1}; \node[right] at (7,0.45) {2};
    \draw[->] (7.4,0.45) -- (8.8,0.45) node[right] {$x$};
    \draw[->] (0,1.2) -- (0,2.4) node[above] {$y$};
    \draw (3.0,0) -- (3.0,0.9); \draw (3.2,0) -- (3.2,0.9); \node[above] at (3.1,0.95) {$dx$};
    \draw[dim] (0.4,0.45) -- (0.4,0.9); \node[right] at (0.45,0.6) {$h$};
    \node[below right] at (0.3,-0.05) {$EI_z = $ constant};
    % temperature profile
    \draw (5.0,0.9) -- (5.0,0); \draw (5.0,0.9) -- (6.0,0.9) -- (5.0,0) -- cycle;
    \node[above] at (6.0,0.95) {$T_u$}; \node[below] at (5.5,-0.05) {$T_l$};
    \draw[dim] (0,-0.9) -- (7,-0.9) node[midway,fill=white,inner sep=1pt] {$L$};
    \node at (3.5,-1.5) {(a)};
    % I section
    \draw (8.9,1.6) -- (9.5,1.6) -- (9.5,1.5) -- (9.25,1.5) -- (9.25,0.9) -- (9.5,0.9) -- (9.5,0.8) -- (8.9,0.8) -- (8.9,0.9) -- (9.15,0.9) -- (9.15,1.5) -- (8.9,1.5) -- cycle;
  \end{scope}
  % ---- (b) deformed beam ----
  \begin{scope}[yshift=-4.6cm]
    \draw (0,0.9) .. controls (2.3,0.0) and (4.7,0.0) .. (7,0.9);
    \draw (0,1.7) .. controls (2.3,0.9) and (4.7,0.9) .. (7,1.7);
    \draw (0,0.9) -- (0,1.7); \draw (7,0.9) -- (7,1.7);
    \fill[hatch] (3.35,0.32) rectangle (3.65,1.12);
    \pic at (0,0.9) {pin}; \pic at (7,0.9) {roller};
    \draw[dim] (3.35,2.1) -- (3.65,2.1); \node[right] at (3.7,2.1) {$dx+\alpha T_u dx$};
    \draw[dim] (3.35,-0.2) -- (3.65,-0.2); \node[right] at (3.7,-0.2) {$dx+\alpha T_l dx$};
    \node at (3.5,-0.9) {(b)};
  \end{scope}
  % ---- (d) fixed-end moments ----
  \begin{scope}[yshift=-8.4cm]
    \draw (0,0) rectangle (7,0.7);
    \pic at (0,0) {pin}; \pic at (7,0) {roller};
    \node[below left] at (0,0) {1}; \node[below right] at (7,0) {2};
    \draw[->,thick] (-0.3,-0.4) arc (-100:120:0.55); \node[above left,align=right] at (-0.2,0.9) {$M^F_{z1} = \dfrac{EI_z\alpha(T_l-T_u)}{h}$};
    \draw[->,thick] (7.3,-0.4) arc (280:60:0.55); \node[above right,align=left] at (7.2,0.9) {$M^F_{z2} = \dfrac{EI_z\alpha(T_l-T_u)}{h}$};
    \node at (3.5,-1.2) {(d)};
  \end{scope}
  % ---- (c) differential element ----
  \begin{scope}[xshift=11.5cm,yshift=-5.5cm]
    \coordinate (O) at (0,5.2);
    \draw (O) -- (-1.0,0); \draw (O) -- (1.0,0);
    \draw[thick] (-0.35,1.6) -- (0.35,1.6) -- (0.6,0.3) -- (-0.6,0.3) -- cycle;
    \draw (-0.2,1.6) rectangle (0.2,1.9);
    \draw[dim] (-0.8,1.6) -- (-0.8,0.3) node[midway,left] {$h$};
    \draw[dim] (-1.4,5.2) -- (-1.4,0.3); \node[left] at (-1.4,2.8) {$\rho$};
    \draw[->] (-0.35,4.3) arc (250:290:1.0); \node[below] at (0,3.9) {$d\theta$};
    \node[right] at (0.4,1.75) {$(1+\alpha T_u)dx$};
    \draw[dim] (-0.6,-0.3) -- (0.6,-0.3); \node[right] at (0.7,-0.3) {$(1+\alpha T_l)dx$};
    \node at (0,-0.9) {(c)};
  \end{scope}
\end{tikzpicture}

\end{fig}
Referring to the above figure, the fixed end moment can be written as $EI\dfrac{\alpha\Delta T}{h}$.
